\chapter{Wstęp}
\indent Duża klasa problemów życia codziennego może być modelowanych przez odpowiadające im problemy grafowe.  
Wiele problemów grafowych jest NP-zupełnych lub NP-trudnych; są też problemy, dla których istnieją algorytmy wielomianowe.  
Dla wielu tych algorytmów kluczowym etapem jest przeprowadzenie wyszukiwania, np. minimalnej drogi.  
Ten etap można znacząco przyśpieszyć wykorzystując kwantowy algorytm wyszukiwania minimum, będący uogólnieniem znanego kwantowego algorytmu wyszukiwania Grovera. \newline
\indent Celem tej pracy jest porównanie klasycznego i kwantowego wariantu algorytmu Dijkstry, służącego do znajdowania najkrótszej drogi w grafie. Algorytm ten zostanie zaimplementowany na trzy różne sposoby, a następnie zostanie sprawdzone, jak te podejścia radzą sobie na różnych rodzajach i rozmiarach grafów.
\newline
\indent W pierwszej części pracy przedstawione zostaną podstawowe zagadnienia teoretyczne, obejmujące definicję grafów, ich rodzaje wykorzystywane w przeprowadzonych badaniach oraz zasadę działania algorytmu Dijkstry. Omówione zostaną również różne sposoby implementacji tego algorytmu.
W dalszej części zaprezentowane zostaną wyniki przeprowadzonych eksperymentów wraz z ich analizą statystyczną, umożliwiającą ocenę rzeczywistej wydajności poszczególnych implementacji w zależności od charakterystyki analizowanych grafów.
Na końcu pracy zamieszczono instrukcję instalacji i konfiguracji środowiska oraz wymaganych bibliotek, umożliwiającą samodzielne uruchomienie aplikacji i odtworzenie przedstawionych eksperymentów.